% GENERATED FILE. Edit canonical lessons, then run npm run generate:books.
% canonical-source-hash: fnv1a64:c20df6a5279b34ca
% canonical-lessons: HI-C250-fail, HI-C250-tijori, HI-C250-mohar, HI-C250-chut, HI-C250-khudra

\chapter{A file, A safe, A stamp, A discount, Retail}
\label{ch:hi-a-file-a-safe-a-stamp-a-discount-retail}
\hlchaptermodality{250}

\begin{chapteropening}
\textbf{By the end of this chapter:} \emph{I can say a file, a safe, a stamp, a discount and retail.}

Complete the last lesson of chapter 250: I can say a file, a safe, a stamp, a discount and retail.
\end{chapteropening}

\section[\texorpdfstring{\hi{फ़ाइल} (fāil)}{फ़ाइल (fāil)}]{\hi{फ़ाइल} (fāil) — a file}

\label{lesson:HI-C250-fail}



\begin{quote}
Before the new one: say the Hindi for a mobile phone, then the Hindi for a document.
\end{quote}

\subsection*{You'll want to know: \hi{फ़ाइल}}
\textbf{\hi{फ़ाइल}} — \emph{fāil} — \textquotedblleft{}a file\textquotedblright{}.

\begin{cousinweb}
\textbf{\hi{फ़ाइल}} (\emph{fāil}) is a loan from English \emph{file}.
\end{cousinweb}

\begin{grammarlens}[title={What you've built}]
One more word for this chapter.
\end{grammarlens}

\subsection*{Your turn}
\begin{itemize}
\raggedright
  \item \emph{Say it:} \emph{fāil}
  \item \emph{Say it:} \emph{fāil}, once more
  \item \emph{Say it:} say \emph{fāil}
  \item \emph{Recall:} say the Hindi for a mobile phone, then the Hindi for a document, then say \emph{fāil} again
\end{itemize}

\begin{tcolorbox}[breakable,colback=teal!4,colframe=teal!35!black,title={Before you move on}]
What does \hi{फ़ाइल} mean? (\textquotedblleft{}A file\textquotedblright{}.) Say it once more.
\end{tcolorbox}
\section[\texorpdfstring{\hi{तिजोरी} (tijorī)}{तिजोरी (tijorī)}]{\hi{तिजोरी} (tijorī) — a safe, a strongbox}

\label{lesson:HI-C250-tijori}



\begin{quote}
Before the new one: say the Hindi for a document, then the Hindi for a file.
\end{quote}

\subsection*{You'll want to know: \hi{तिजोरी}}
\textbf{\hi{तिजोरी}} — \emph{tijorī} — \textquotedblleft{}a safe, a strongbox\textquotedblright{}.

\begin{cousinweb}
The origin of \textbf{\hi{तिजोरी}} (\emph{tijorī}) is uncertain. It is sometimes traced to Portuguese \emph{tesouro}, \textquotedblleft{}treasure, treasury\textquotedblright{}, but no source settles it, and CDIAL has no entry for it.
\end{cousinweb}

\begin{grammarlens}[title={What you've built}]
One more word for this chapter.
\end{grammarlens}

\subsection*{Your turn}
\begin{itemize}
\raggedright
  \item \emph{Say it:} \emph{tijorī}
  \item \emph{Say it:} \emph{tijorī}, once more
  \item \emph{Say it:} say \emph{tijorī}
  \item \emph{Recall:} say the Hindi for a document, then the Hindi for a file, then say \emph{tijorī} again
\end{itemize}

\begin{tcolorbox}[breakable,colback=teal!4,colframe=teal!35!black,title={Before you move on}]
What does \hi{तिजोरी} mean? (\textquotedblleft{}A safe, a strongbox\textquotedblright{}.) Say it once more.
\end{tcolorbox}
\section[\texorpdfstring{\hi{मोहर} (mohar)}{मोहर (mohar)}]{\hi{मोहर} (mohar) — a stamp, a seal}

\label{lesson:HI-C250-mohar}



\begin{quote}
Before the new one: say the Hindi for a file, then the Hindi for a safe.
\end{quote}

\subsection*{You'll want to know: \hi{मोहर}}
\textbf{\hi{मोहर}} — \emph{mohar} — \textquotedblleft{}a stamp, a seal\textquotedblright{}.

\begin{cousinweb}
\textbf{\hi{मोहर}} (\emph{mohar}) is a loan, usually traced to Persian \emph{muhr}, \textquotedblleft{}seal, stamp\textquotedblright{}.
\end{cousinweb}

\begin{grammarlens}[title={What you've built}]
One more word for this chapter.
\end{grammarlens}

\subsection*{Your turn}
\begin{itemize}
\raggedright
  \item \emph{Say it:} \emph{mohar}
  \item \emph{Say it:} \emph{mohar}, once more
  \item \emph{Say it:} say \emph{mohar}
  \item \emph{Recall:} say the Hindi for a file, then the Hindi for a safe, then say \emph{mohar} again
\end{itemize}

\begin{tcolorbox}[breakable,colback=teal!4,colframe=teal!35!black,title={Before you move on}]
What does \hi{मोहर} mean? (\textquotedblleft{}A stamp, a seal\textquotedblright{}.) Say it once more.
\end{tcolorbox}
\section[\texorpdfstring{\hi{छूट} (chūṭ)}{छूट (chūṭ)}]{\hi{छूट} (chūṭ) — a discount}

\label{lesson:HI-C250-chut}



\begin{quote}
Before the new one: say the Hindi for a safe, then the Hindi for a stamp.
\end{quote}

\subsection*{You'll want to know: \hi{छूट}}
\textbf{\hi{छूट}} — \emph{chūṭ} — \textquotedblleft{}a discount\textquotedblright{}.

\begin{cousinweb}
\textbf{\hi{छूट}} (\emph{chūṭ}) is a tadbhava, a word inherited and reshaped by sound change. It comes from \emph{*kṣuṭyate}, \textquotedblleft{}is released\textquotedblright{}, a form CDIAL reconstructs from its descendants (CDIAL 3707). \emph{chūṭ}, \textquotedblleft{}a let-off\textquotedblright{}, is the noun of \emph{chūṭnā}, \textquotedblleft{}to be released\textquotedblright{}.
\end{cousinweb}

\begin{grammarlens}[title={What you've built}]
One more word for this chapter.
\end{grammarlens}

\subsection*{Your turn}
\begin{itemize}
\raggedright
  \item \emph{Say it:} \emph{chūṭ}
  \item \emph{Say it:} \emph{chūṭ}, once more
  \item \emph{Say it:} say \emph{chūṭ}
  \item \emph{Recall:} say the Hindi for a safe, then the Hindi for a stamp, then say \emph{chūṭ} again
\end{itemize}

\begin{tcolorbox}[breakable,colback=teal!4,colframe=teal!35!black,title={Before you move on}]
What does \hi{छूट} mean? (\textquotedblleft{}A discount\textquotedblright{}.) Say it once more.
\end{tcolorbox}
\section[\texorpdfstring{\hi{खुदरा} (khuḍrā)}{खुदरा (khuḍrā)}]{\hi{खुदरा} (khuḍrā) — retail, small change}

\label{lesson:HI-C250-khudra}



\begin{quote}
Before the new one: say the Hindi for a stamp, then the Hindi for a discount.
\end{quote}

\subsection*{You'll want to know: \hi{खुदरा}}
\textbf{\hi{खुदरा}} — \emph{khuḍrā} — \textquotedblleft{}retail, small change\textquotedblright{}.

\begin{cousinweb}
\textbf{\hi{खुदरा}} (\emph{khudrā}) is a loan, usually traced to Persian \emph{khurda}, \textquotedblleft{}small, small change\textquotedblright{}.
\end{cousinweb}

\begin{grammarlens}[title={What you've built}]
One more word for this chapter.
\end{grammarlens}

\subsection*{Your turn}
\begin{itemize}
\raggedright
  \item \emph{Say it:} \emph{khuḍrā}
  \item \emph{Say it:} \emph{khuḍrā}, once more
  \item \emph{Say it:} say \emph{khuḍrā}
  \item \emph{Recall:} say the Hindi for a stamp, then the Hindi for a discount, then say \emph{khuḍrā} again
\end{itemize}

\begin{tcolorbox}[breakable,colback=teal!4,colframe=teal!35!black,title={Before you move on}]
What does \hi{खुदरा} mean? (\textquotedblleft{}Retail, small change\textquotedblright{}.) Say it once more.
\end{tcolorbox}
